% Drafted against audit/outline.md §8. Two additions to the outline's five
% bullets, both of which became limitations only after the disclosed-settings
% rerun: the undisclosed decay rate, and the fact that the disclosed window
% leaves no room to appeal to a longer one. Recorded in the outline.
%
% No new evidence. Everything here is carried from §3, §4 and §7.

\section{Limitations}\label{sec:limitations}

Everything in this section constrains how the results should be read. The two
that matter most run in opposite directions, and the paper does not claim they
cancel.

\subsection{Seventeen factors against roughly seventy thousand}

This is the largest single difference between the engine here and Goldman's,
and it is not fixable from public data. More independent sources of risk means
more offsetting available in the joint tail, so a seventeen-factor book should
show \emph{less} diversification than a book with a richer factor set. The
direction of that bias on the level is therefore known, and it makes the level
agreement in Section~\ref{sec:results} slightly surprising rather than
suspicious.

The direction on the steadiness is not known, and that is the more important
gap. A larger factor set might average out regime shifts in any one
correlation, which would push toward the disclosed stability; or it might not,
if the shifts are common across factors, which is what a macro regime change
usually means. Nothing in public data settles it. Section~\ref{sec:assumptions}
lists it as a surviving candidate and declines to guess, and the central result
should be read with that candidate open.

\subsection{Non-synchronous closes}

Tokyo, London and New York do not close at the same instant. A move that
happens in New York after the Tokyo close appears in the Japanese series only
the following day, which mechanically depresses measured cross-market
correlation and therefore \emph{overstates} diversification.

This runs opposite to the factor-count bias. Neither is quantified here, and
the two are not assumed to offset: they are stated so that a reader knows the
measured level sits between two unestimated biases of unknown relative size.
This is also why neither is used anywhere to argue that the result would move
in a particular direction if corrected.

\subsection{What the calendar intersection costs}

Restricting the panel to days on which every market traded discards about 9\%
of trading days and leaves 225 to 233 observations a year against the 250 to
253 Goldman reports. The alternative --- forward-filling closed markets ---
would plant runs of exact zeros in the scenario set and understate those
factors' risk, which biases the quantity being measured rather than merely
reducing the precision of the measurement. The smaller sample is the better of
two imperfect options, but it does cost power in every test in
Section~\ref{sec:validation}, and the annual figures rest on fewer observations
than the disclosures they are compared against.

A related quantisation affects the rate factors. The Treasury curve is quoted
to the basis point, so at near-zero policy rates the two-year yield shows no
change at all on 21.1\% of days across the panel and 43.1\% of days in the
second half of 2020. That is a real feature of the quoted series rather than a
data error, and it slightly understates true rate variability through the
near-zero years.

\subsection{Instrument substitutions}

Commodity exposure runs through exchange-traded funds rather than front-month
futures, for the reason given in Section~\ref{sec:sample}. Funds carry roll
yield, which drags them away from spot over long horizons. Over the one-day
horizon the engine uses it is second order, but it is a real difference from a
spot book and it is not corrected.

The rate sleeve mixes US par coupon bonds with euro zero-coupon bonds, each in
its source curve's native convention. A zero has longer duration than a par
bond of the same maturity, so the two are not interchangeable per unit of
notional. Sizing targets risk rather than notional, which absorbs the
difference for the purpose of matching the disclosed split, but the instrument
sets are not identical and a reader comparing notionals across sleeves in
Table~\ref{tab:book} should not read them as comparable exposures.

\subsection{What the recalibrated book cannot be used for}

The recalibrated book is sized each year using figures Goldman published after
that year ended: the 2022 sizing uses a figure released in February 2023. That
is appropriate for the structural comparison it is used for, because the
question is whether a book with the disclosed risk split generates the
disclosed co-movement, and the split is the input. It would be indefensible as
a forecast. This is why only the fixed book, sized once on 2019 and never
touched again, is backtested, and why the 2020--2025 rows of
Table~\ref{tab:backtest} are out of sample with respect to sizing.

\subsection{Two limits the disclosed settings impose}

Adopting Goldman's stated method closed some latitude and opened a different
constraint.

The decay rate behind the disclosed age weighting is not published. The filings
say the importance of the data reduces over time without saying how fast, so
the scheme is tested in kind but not at Goldman's own rate. The two rates run
here were chosen for their half-lives and both move the steadiness the same
way, which is what makes the conclusion in Section~\ref{sec:assumptions} robust
to not knowing the rate, but it remains a conclusion about a family of decay
rates rather than about Goldman's.

The window leaves no further room. Five years is what the filings disclose, so
the monotone improvement in steadiness with window length cannot be extended by
appeal to a longer one. That is a limitation on the argument's escape routes
rather than on its evidence: the result is bounded by what the firm says it
does.

\subsection{What validation does not cover}

Finally, and most importantly for interpretation: Section~\ref{sec:validation}
establishes that the engine computes what Section~\ref{sec:method} says it
computes. It does not establish --- and no test on public data could --- that
the proxy book resembles Goldman's trading book. The level agreement is
evidence about the proxy, and the steadiness gap is evidence about what a
static proxy cannot do. Neither is a measurement of anything inside the firm.
